\section{数理逻辑}

\subsection{数理逻辑概览}

离散数学中的数理逻辑主要包括命题逻辑和谓词逻辑。命题逻辑研究命题之间如何通过联结词组合、等值变形和推理；谓词逻辑进一步分析命题内部的个体、谓词和量词，从而表达“所有”“存在”等更细的逻辑结构。

\subsection{命题逻辑}

\subsubsection{命题与真值}

\begin{enumerate}
	\item \textbf{命题}：能够判断真假的陈述句。疑问句、祈使句、感叹句通常不是命题。
	\item \textbf{真值}：命题的真假值，常用 \(1\) 表示真，\(0\) 表示假。
	\item \textbf{原子命题}：不能再分解为更简单命题的命题。
	\item \textbf{复合命题}：由原子命题和联结词组成的命题。
	\item 真值判断常受到时间、区域和判定标准影响，但只要客观上能判定真假，就可作为命题。
\end{enumerate}

\begin{question}{例题1}
	判断下列语句是否为命题。
	\begin{enumerate}
		\item 北京是中国的首都。
		\item 请关门！
		\item \(2x+4\ge 10\)。
		\item 地球外的星球上有生命。
	\end{enumerate}
	
	\begin{solution}
	(1) 是命题，且为真命题；(2) 不是命题，因为它是祈使句；(3) 不是命题，因为变量 \(x\) 未指定，真假未定；(4) 是命题，虽然目前可能无法判定真值，但它客观上或真或假。
	\end{solution}
\end{question}

\subsubsection{命题联结词}

\begin{center}
	\begin{tabular}{cccc}
		\toprule
		联结词 & 符号 & 英文含义 & 读法\\
		\midrule
		否定 & \(\neg\) & NOT & 非 \(p\)\\
		合取 & \(\land\) & AND & \(p\) 并且 \(q\)\\
		析取 & \(\lor\) & OR & \(p\) 或 \(q\)\\
		蕴含 & \(\to\) & if--then & 如果 \(p\)，则 \(q\)\\
		等价 & \(\leftrightarrow\) & if and only if & \(p\) 当且仅当 \(q\)\\
		\bottomrule
	\end{tabular}
\end{center}

运算优先级为
\[
\neg\;>\;\land\;>\;\lor\;>\;\to\;>\;\leftrightarrow.
\]
例如
\[
\neg p\to q\lor r\quad \text{应理解为}\quad (\neg p)\to(q\lor r).
\]

\subsubsection{常见真值表}

\begin{center}
	\begin{tabular}{cccccc}
		\toprule
		\(p\) & \(q\) & \(p\land q\) & \(p\lor q\) & \(p\to q\) & \(p\leftrightarrow q\)\\
		\midrule
		0&0&0&0&1&1\\
		0&1&0&1&1&0\\
		1&0&0&1&0&0\\
		1&1&1&1&1&1\\
		\bottomrule
	\end{tabular}
\end{center}

\begin{enumerate}
	\item 蕴含式 \(p\to q\) 只有在 \(p=1,q=0\) 时为假，其余均为真。
	\item 等价式 \(p\leftrightarrow q\) 在 \(p,q\) 真值相同时为真，不同时为假。
	\item 自然语言中的“或”有时表示相容析取，有时表示不可兼析取，符号化时要结合语义判断。
\end{enumerate}

\subsubsection{不可兼或、与非和或非}

\begin{enumerate}
	\item \textbf{不可兼或}又称异或，记作 \(p\veebar q\) 或 \(p\oplus q\)，当且仅当 \(p,q\) 真值不同时为真：
	\[
	p\oplus q\Leftrightarrow (p\lor q)\land\neg(p\land q).
	\]
	\[
	p\oplus q\Leftrightarrow \neg(p\leftrightarrow q).
	\]
	\item \textbf{与非}：
	\[
	p\uparrow q\Leftrightarrow \neg(p\land q).
	\]
	当且仅当 \(p,q\) 同真时，\(p\uparrow q\) 为假。
	\item \textbf{或非}：
	\[
	p\downarrow q\Leftrightarrow \neg(p\lor q).
	\]
	当且仅当 \(p,q\) 同假时，\(p\downarrow q\) 为真。
\end{enumerate}

\begin{center}
	\begin{tabular}{ccccc}
		\toprule
		\(p\)&\(q\)&\(p\oplus q\)&\(p\uparrow q\)&\(p\downarrow q\)\\
		\midrule
		0&0&0&1&1\\
		0&1&1&1&0\\
		1&0&1&1&0\\
		1&1&0&0&0\\
		\bottomrule
	\end{tabular}
\end{center}

\subsubsection{命题公式分类}

\begin{enumerate}
	\item \textbf{重言式}：任意赋值下均为真，也称永真式。
	\item \textbf{矛盾式}：任意赋值下均为假，也称永假式。
	\item \textbf{可满足式}：至少存在一个成真赋值。
	\item \textbf{偶然式}：既有成真赋值，又有成假赋值。
\end{enumerate}

\subsection{命题等值演算}

\subsubsection{等值式}

若公式 \(A\leftrightarrow B\) 是重言式，则称 \(A\) 与 \(B\) 等值，记作
\[
A\Leftrightarrow B.
\]
等值演算的用途包括：判定公式类型、证明公式等值、化简复杂命题公式、求范式。

\subsubsection{常用等值式}

\begin{center}
	\begin{tabular}{cl}
		\toprule
		名称 & 公式\\
		\midrule
		双否律 & \(\neg\neg A\Leftrightarrow A\)\\
		幂等律 & \(A\land A\Leftrightarrow A,\quad A\lor A\Leftrightarrow A\)\\
		交换律 & \(A\land B\Leftrightarrow B\land A,\quad A\lor B\Leftrightarrow B\lor A\)\\
		结合律 & \((A\land B)\land C\Leftrightarrow A\land(B\land C)\)\\
		       & \((A\lor B)\lor C\Leftrightarrow A\lor(B\lor C)\)\\
		分配律 & \(A\land(B\lor C)\Leftrightarrow(A\land B)\lor(A\land C)\)\\
		       & \(A\lor(B\land C)\Leftrightarrow(A\lor B)\land(A\lor C)\)\\
		德摩根律 & \(\neg(A\land B)\Leftrightarrow\neg A\lor\neg B\)\\
		          & \(\neg(A\lor B)\Leftrightarrow\neg A\land\neg B\)\\
		吸收律 & \(A\land(A\lor B)\Leftrightarrow A,\quad A\lor(A\land B)\Leftrightarrow A\)\\
		同一律 & \(A\land 1\Leftrightarrow A,\quad A\lor 0\Leftrightarrow A\)\\
		零律 & \(A\land 0\Leftrightarrow0,\quad A\lor1\Leftrightarrow1\)\\
		排中律 & \(A\lor\neg A\Leftrightarrow1\)\\
		矛盾律 & \(A\land\neg A\Leftrightarrow0\)\\
		蕴含式 & \(A\to B\Leftrightarrow\neg A\lor B\)\\
		蕴含否定 & \(\neg(A\to B)\Leftrightarrow A\land\neg B\)\\
		等价式 & \(A\leftrightarrow B\Leftrightarrow(A\to B)\land(B\to A)\)\\
		等价式 & \(A\leftrightarrow B\Leftrightarrow(A\land B)\lor(\neg A\land\neg B)\)\\
		假言易位 & \(A\to B\Leftrightarrow\neg B\to\neg A\)\\
		归谬式 & \(A\to B\Leftrightarrow(A\land\neg B)\to0\)\\
		\bottomrule
	\end{tabular}
\end{center}

\begin{question}{例题2}
	证明
	\[
	q\to(p\to r)\Leftrightarrow(p\land q)\to r.
	\]
	
	\begin{solution}
	\begin{align*}
	q\to(p\to r)
	&\Leftrightarrow \neg q\lor(\neg p\lor r)\\
	&\Leftrightarrow \neg p\lor\neg q\lor r\\
	&\Leftrightarrow \neg(p\land q)\lor r\\
	&\Leftrightarrow (p\land q)\to r.
	\end{align*}
	\end{solution}
\end{question}

\subsection{命题蕴含与推理规则}

\subsubsection{永真蕴含}

若 \(A\to B\) 是重言式，则称 \(A\) 永真蕴含 \(B\)，记作
\[
A\Rightarrow B.
\]
注意：\(A\Rightarrow B\) 是元语言中的推理关系；\(A\to B\) 是命题逻辑内部的公式。

\subsubsection{常用蕴含式}

\begin{center}
	\begin{tabular}{cl}
		\toprule
		名称 & 蕴含式\\
		\midrule
		化简律 & \(A\land B\Rightarrow A,\quad A\land B\Rightarrow B\)\\
		附加律 & \(A\Rightarrow A\lor B,\quad B\Rightarrow A\lor B\)\\
		假言推理 & \(A,\ A\to B\Rightarrow B\)\\
		拒取式 & \(A\to B,\ \neg B\Rightarrow \neg A\)\\
		析取三段论 & \(A\lor B,\ \neg A\Rightarrow B\)\\
		合取引入 & \(A,\ B\Rightarrow A\land B\)\\
		假言三段论 & \(A\to B,\ B\to C\Rightarrow A\to C\)\\
		等价三段论 & \(A\leftrightarrow B,\ B\leftrightarrow C\Rightarrow A\leftrightarrow C\)\\
		构造性二难 & \(A\to B,\ C\to D,\ A\lor C\Rightarrow B\lor D\)\\
		归结式 & \(A\lor B,\ \neg A\lor C\Rightarrow B\lor C\)\\
		\bottomrule
	\end{tabular}
\end{center}

\subsubsection{推理规则与证明方法}

\begin{enumerate}
	\item \textbf{P 规则}：推导的任意步骤可引入前提。
	\item \textbf{T 规则}：若前面公式永真蕴含公式 \(S\)，则可引入 \(S\)。
	\item \textbf{CP 规则}：若
	\[
	H_1\land H_2\land\cdots\land H_n\land R\Rightarrow S,
	\]
	则
	\[
	H_1\land H_2\land\cdots\land H_n\Rightarrow R\to S.
	\]
	\item \textbf{直接证法}：从前提出发，用推理规则直接推出结论。
	\item \textbf{附加前提证明法}：若结论形如 \(R\to S\)，可附加前提 \(R\)，推出 \(S\)。
	\item \textbf{反证法}：要证 \(H_1,\dots,H_n\Rightarrow C\)，只需证明
	\[
	H_1\land\cdots\land H_n\land\neg C
	\]
	为矛盾式。
\end{enumerate}

\begin{question}{例题3}
	证明
	\[
	(p\lor q)\land(p\leftrightarrow r)\land(q\to s)\Rightarrow s\lor r.
	\]
	
	\begin{solution}
	若结论 \(s\lor r\) 为假，则
	\[
	\neg(s\lor r)\Leftrightarrow\neg s\land\neg r.
	\]
	由 \(p\leftrightarrow r\) 得 \(p\to r\)。再由 \(\neg r\)，用拒取式得 \(\neg p\)。由 \(q\to s\) 和 \(\neg s\)，用拒取式得 \(\neg q\)。于是有 \(\neg p\land\neg q\)，这与前提 \(p\lor q\) 矛盾。因此结论 \(s\lor r\) 必成立。
	\end{solution}
\end{question}

\begin{question}{例题4}
	用附加前提法证明：
	\[
	(p\land q)\to r,\quad \neg s\lor p,\quad q\vdash s\to r.
	\]
	
	\begin{solution}
	\begin{center}
	\begin{tabular}{cll}
		\toprule
		序号 & 公式 & 理由\\
		\midrule
		1&\((p\land q)\to r\)&前提\\
		2&\(\neg s\lor p\)&前提\\
		3&\(q\)&前提\\
		4&\(s\)&附加前提\\
		5&\(p\)&2,4 析取三段论\\
		6&\(p\land q\)&5,3 合取引入\\
		7&\(r\)&1,6 假言推理\\
		8&\(s\to r\)&4--7 附加前提法\\
		\bottomrule
	\end{tabular}
	\end{center}
	\end{solution}
\end{question}

\subsection{对偶与范式}

\subsubsection{对偶式}

在命题公式 \(A\) 中，把联结词 \(\land\) 与 \(\lor\) 互换，把常量 \(1\) 与 \(0\) 互换，所得公式称为 \(A\) 的对偶式，记作 \(A^*\)。例如
\[
(p\lor q)\land0
\]
的对偶式为
\[
(p\land q)\lor1.
\]

\textbf{对偶原则}：若 \(A\Leftrightarrow B\)，则 \(A^*\Leftrightarrow B^*\)。

\subsubsection{析取范式与合取范式}

\begin{enumerate}
	\item \(p\) 和 \(\neg p\) 称为\textbf{文字}。
	\item 有限个文字的析取称为\textbf{析取式}；有限个文字的合取称为\textbf{合取式}。
	\item 有限个合取式的析取称为\textbf{析取范式}。
	\item 有限个析取式的合取称为\textbf{合取范式}。
	\item 任意命题公式都存在与其等值的析取范式和合取范式。
\end{enumerate}

\subsubsection{主析取范式与主合取范式}

含 \(n\) 个命题变元的合取式中，每个变元或其否定出现且仅出现一次，称为\textbf{小项}。小项只在一个赋值下为真。全体成真小项的析取就是主析取范式。

含 \(n\) 个命题变元的析取式中，每个变元或其否定出现且仅出现一次，称为\textbf{大项}。大项只在一个赋值下为假。全体成假大项的合取就是主合取范式。

\begin{center}
	\begin{tabular}{ccc}
		\toprule
		小项 & 成真赋值 & 名称\\
		\midrule
		\(\neg p\land\neg q\)&00&\(m_0\)\\
		\(\neg p\land q\)&01&\(m_1\)\\
		\(p\land\neg q\)&10&\(m_2\)\\
		\(p\land q\)&11&\(m_3\)\\
		\bottomrule
	\end{tabular}
	\qquad
	\begin{tabular}{ccc}
		\toprule
		大项 & 成假赋值 & 名称\\
		\midrule
		\(p\lor q\)&00&\(M_0\)\\
		\(p\lor\neg q\)&01&\(M_1\)\\
		\(\neg p\lor q\)&10&\(M_2\)\\
		\(\neg p\lor\neg q\)&11&\(M_3\)\\
		\bottomrule
	\end{tabular}
\end{center}

\begin{question}{例题5}
	求 \(P\to Q\) 的主析取范式和主合取范式。
	
	\begin{solution}
	\[
	P\to Q\Leftrightarrow\neg P\lor Q.
	\]
	其真值为真时的赋值为 \(00,01,11\)，故主析取范式为
	\[
	(\neg P\land\neg Q)\lor(\neg P\land Q)\lor(P\land Q).
	\]
	其唯一成假赋值为 \(10\)，对应大项 \(\neg P\lor Q\)，故主合取范式为
	\[
	\neg P\lor Q.
	\]
	\end{solution}
\end{question}

\subsubsection{联结词完备集}

若任意命题公式都可只用联结词集合 \(S\) 中的联结词表示，则 \(S\) 称为联结词完备集。常见完备集包括：
\[
\{\neg,\land,\lor\},\quad
\{\neg,\land\},\quad
\{\neg,\lor\},\quad
\{\neg,\to\},\quad
\{\uparrow\},\quad
\{\downarrow\}.
\]

例如
\[
P\to Q\Leftrightarrow \neg P\lor Q\Leftrightarrow \neg(P\land\neg Q),
\]
\[
P\to Q\Leftrightarrow P\uparrow(Q\uparrow Q),
\]
\[
P\to Q\Leftrightarrow ((P\downarrow P)\downarrow Q)\downarrow((P\downarrow P)\downarrow Q).
\]

\subsection{谓词逻辑}

\subsubsection{基本概念}

\begin{enumerate}
	\item \textbf{个体}：客观世界中可独立存在的具体或抽象对象。
	\item \textbf{个体常量}：表示特定个体的符号，如 \(a,b,c\)。
	\item \textbf{个体变量}：泛指不确定个体的符号，如 \(x,y,z\)。
	\item \textbf{谓词}：表示个体性质或个体间关系的符号，如 \(P(x)\)、\(R(x,y)\)。
	\item \textbf{个体域}：个体变量的取值范围，也称论域。
	\item \textbf{谓词函数}：由谓词和个体变元组成的表达式。谓词函数本身未必是命题，只有个体变元确定或被量词约束后才成为命题。
\end{enumerate}

\subsubsection{量词}

\begin{enumerate}
	\item \textbf{全称量词} \(\forall\)：\(\forall xP(x)\) 读作“对所有 \(x\)，\(P(x)\) 成立”。
	\item \textbf{存在量词} \(\exists\)：\(\exists xP(x)\) 读作“存在 \(x\)，使 \(P(x)\) 成立”。
	\item \textbf{辖域}：量词作用的公式范围。
	\item \textbf{约束变元}：在同名量词辖域内出现的变量。
	\item \textbf{自由变元}：不受同名量词约束的变量。
\end{enumerate}

\begin{question}{例题6}
	设 \(L(x)\)：\(x\) 是人，\(R(x,y)\)：\(x\) 认识 \(y\)，\(E(x,y)\)：\(x\) 等于 \(y\)。符号化“有人认识所有不是自己的人”。
	
	\begin{solution}
	\[
	\exists x\Bigl(L(x)\land \forall y\bigl((L(y)\land\neg E(x,y))\to R(x,y)\bigr)\Bigr).
	\]
	\end{solution}
\end{question}

\subsubsection{前束范式}

若一个谓词公式中所有量词均位于公式最前端，并且其辖域延伸到整个后续公式，则称该公式为\textbf{前束范式}。例如
\[
(\forall x)(\exists y)(P(x)\to Q(y))
\]
是前束范式。

\begin{question}{例题7}
	把
	\[
	(\forall x)P(x)\to(\exists x)Q(x)
	\]
	转化为前束范式。
	
	\begin{solution}
	为避免变量混淆，先将右侧变量改名：
	\[
	(\forall x)P(x)\to(\exists y)Q(y).
	\]
	然后消去蕴含并移动量词：
	\begin{align*}
	(\forall x)P(x)\to(\exists y)Q(y)
	&\Leftrightarrow \neg(\forall x)P(x)\lor(\exists y)Q(y)\\
	&\Leftrightarrow (\exists x)\neg P(x)\lor(\exists y)Q(y)\\
	&\Leftrightarrow (\exists x)(\exists y)(\neg P(x)\lor Q(y))\\
	&\Leftrightarrow (\exists x)(\exists y)(P(x)\to Q(y)).
	\end{align*}
	\end{solution}
\end{question}

\subsubsection{谓词演算常用等价式}

\begin{enumerate}
	\item 命题逻辑等价式可作为谓词公式的代换实例：
	\[
	P(x)\to Q(x)\Leftrightarrow\neg P(x)\lor Q(x).
	\]
	\item 量词否定：
	\[
	\neg(\forall x)P(x)\Leftrightarrow(\exists x)\neg P(x),
	\]
	\[
	\neg(\exists x)P(x)\Leftrightarrow(\forall x)\neg P(x).
	\]
	\item 量词分配：
	\[
	(\forall x)(P(x)\land Q(x))\Leftrightarrow(\forall x)P(x)\land(\forall x)Q(x),
	\]
	\[
	(\exists x)(P(x)\lor Q(x))\Leftrightarrow(\exists x)P(x)\lor(\exists x)Q(x).
	\]
	\item 单向分配：
	\[
	(\forall x)P(x)\lor(\forall x)Q(x)\Rightarrow(\forall x)(P(x)\lor Q(x)),
	\]
	\[
	(\exists x)(P(x)\land Q(x))\Rightarrow(\exists x)P(x)\land(\exists x)Q(x).
	\]
	\item 辖域扩张与收缩。若 \(x\) 不在 \(Q\) 中自由出现，则
	\[
	(\forall x)(P(x)\land Q)\Leftrightarrow(\forall x)P(x)\land Q,
	\]
	\[
	(\exists x)(P(x)\lor Q)\Leftrightarrow(\exists x)P(x)\lor Q.
	\]
	\item 量词与蕴含。若 \(x\) 不在 \(Q\) 中自由出现，则
	\[
	(\forall x)(P(x)\to Q)\Leftrightarrow(\exists x)P(x)\to Q,
	\]
	\[
	(\exists x)(P(x)\to Q)\Leftrightarrow(\forall x)P(x)\to Q,
	\]
	\[
	(\forall x)(Q\to P(x))\Leftrightarrow Q\to(\forall x)P(x),
	\]
	\[
	(\exists x)(Q\to P(x))\Leftrightarrow Q\to(\exists x)P(x).
	\]
\end{enumerate}

\subsubsection{有限个体域中消去量词}

若个体域 \(D=\{a_1,a_2,\dots,a_n\}\)，则
\[
(\forall x)P(x)\Leftrightarrow P(a_1)\land P(a_2)\land\cdots\land P(a_n),
\]
\[
(\exists x)P(x)\Leftrightarrow P(a_1)\lor P(a_2)\lor\cdots\lor P(a_n).
\]

\begin{question}{例题8}
	设个体域 \(D=\{a,b,c\}\)，消去公式
	\[
	(\exists x)F(x)\to(\forall y)G(y)
	\]
	中的量词。
	
	\begin{solution}
	\[
	(\exists x)F(x)\to(\forall y)G(y)
	\Leftrightarrow
	(F(a)\lor F(b)\lor F(c))\to(G(a)\land G(b)\land G(c)).
	\]
	\end{solution}
\end{question}

\subsubsection{谓词演算推理规则}

网页图片中的四条量词推理规则整理如下：

\begin{center}
	\begin{tabular}{ccp{8cm}}
		\toprule
		规则 & 形式 & 使用条件\\
		\midrule
		全称指定 US & \((\forall x)P(x)\Rightarrow P(c)\) & \(c\) 代表个体域中的任意元素\\
		存在指定 ES & \((\exists x)P(x)\Rightarrow P(c)\) & \(c\) 代表某个满足条件的个体，每次使用要引入新的个体常量\\
		全称推广 UG & \(P(c)\Rightarrow(\forall x)P(x)\) & \(c\) 必须能代表个体域中的任意元素\\
		存在推广 EG & \(P(c)\Rightarrow(\exists x)P(x)\) & \(c\) 是个体域中的某个元素即可\\
		\bottomrule
	\end{tabular}
\end{center}

\begin{question}{例题9}
	证明：
	\[
	(\forall x)(M(x)\to D(x)),\quad(\exists x)(M(x)\land S(x))\vdash(\exists x)(D(x)\land S(x)).
	\]
	
	\begin{solution}
	\begin{center}
	\begin{tabular}{cll}
		\toprule
		序号 & 公式 & 理由\\
		\midrule
		1&\((\forall x)(M(x)\to D(x))\)&前提\\
		2&\((\exists x)(M(x)\land S(x))\)&前提\\
		3&\(M(a)\land S(a)\)&2, ES\\
		4&\(M(a)\)&3, 化简\\
		5&\(S(a)\)&3, 化简\\
		6&\(M(a)\to D(a)\)&1, US\\
		7&\(D(a)\)&4,6, 假言推理\\
		8&\(D(a)\land S(a)\)&7,5, 合取引入\\
		9&\((\exists x)(D(x)\land S(x))\)&8, EG\\
		\bottomrule
	\end{tabular}
	\end{center}
	\end{solution}
\end{question}

\subsection{逻辑部分考试补充}

\subsubsection{合式公式判定}

命题逻辑中，合式公式按以下规则递归定义：
\begin{enumerate}
	\item 单个命题变元是合式公式。
	\item 若 \(A\) 是合式公式，则 \(\neg A\) 是合式公式。
	\item 若 \(A,B\) 是合式公式，则 \((A\land B),(A\lor B),(A\to B),(A\leftrightarrow B)\) 是合式公式。
	\item 只由上述规则有限次生成的符号串才是合式公式。
\end{enumerate}

考试中要区分对象语言符号和元语言符号：\(\to,\leftrightarrow\) 是公式内部联结词；\(\Rightarrow,\Leftrightarrow\) 常用于表示推理关系和等值关系，一般不作为合式公式中的联结词。

\subsubsection{赋值个数与小项大项数量}

若命题公式含 \(n\) 个不同命题变元，则共有
\[
2^n
\]
种赋值。含 \(n\) 个命题变元的小项共有 \(2^n\) 个，大项也共有 \(2^n\) 个。

\begin{enumerate}
	\item 全体小项的析取为永真式：
	\[
	\bigvee_{i=0}^{2^n-1}m_i\Leftrightarrow 1.
	\]
	\item 全体小项的合取为矛盾式：
	\[
	\bigwedge_{i=0}^{2^n-1}m_i\Leftrightarrow 0.
	\]
	\item 全体大项的合取为矛盾式：
	\[
	\bigwedge_{i=0}^{2^n-1}M_i\Leftrightarrow 0.
	\]
	\item 全体大项的析取为永真式：
	\[
	\bigvee_{i=0}^{2^n-1}M_i\Leftrightarrow 1.
	\]
\end{enumerate}

\subsubsection{常见自然语言符号化}

设 \(P\)：我生病，\(Q\)：我去学校。
\begin{center}
	\begin{tabular}{p{7cm}c}
		\toprule
		自然语言 & 符号化\\
		\midrule
		若我生病，则我不去学校 & \(P\to\neg Q\)\\
		只有在生病时，我才不去学校 & \(\neg Q\to P\)\\
		当且仅当我生病时，我才不去学校 & \(\neg Q\leftrightarrow P\)\\
		若我不生病，则我一定去学校 & \(\neg P\to Q\)\\
		除非你努力，否则你将失败 & \(\neg P\to Q\)，等价于 \(\neg Q\to P\)\\
		虽然你努力了，但还是失败了 & \(P\land Q\)\\
		\bottomrule
	\end{tabular}
\end{center}

\subsubsection{前束范式步骤}

把谓词公式化为前束范式常按以下步骤：
\begin{enumerate}
	\item 消去 \(\to,\leftrightarrow\)。
	\item 将否定深入到原子谓词前，使用量词否定律。
	\item 对约束变元换名，避免不同量词使用同名变元，也避免与自由变元冲突。
	\item 利用辖域扩张和收缩，把量词逐步移到公式最前端。
\end{enumerate}

\begin{question}{例题10}
	求
	\[
	\forall xF(x)\land\neg(\exists xG(x))
	\]
	的前束范式。
	
	\begin{solution}
	先将第二个量词换名并否定深入：
	\[
	\forall xF(x)\land\neg(\exists yG(y))
	\Leftrightarrow
	\forall xF(x)\land\forall y\neg G(y).
	\]
	于是前束范式为
	\[
	\forall x\forall y(F(x)\land\neg G(y)).
	\]
	\end{solution}
\end{question}

\subsubsection{自由变元与约束变元}

判断自由变元和约束变元时，只看该变量的每一次出现是否落在同名量词的辖域内。同一个字母可以既有自由出现又有约束出现。

\begin{question}{例题11}
	指出公式
	\[
	\forall x(A(x)\to B(y,x))\land\exists zC(y,z)\to D(x)
	\]
	中的自由变元和约束变元。
	
	\begin{solution}
	在 \(\forall x(A(x)\to B(y,x))\) 中，两个 \(x\) 是约束出现，\(y\) 是自由出现。在 \(\exists zC(y,z)\) 中，\(z\) 是约束出现，\(y\) 是自由出现。最后 \(D(x)\) 中的 \(x\) 不在前面 \(\forall x\) 的辖域内，所以是自由出现。
	
	因此自由变元为 \(x,y\)，约束变元为 \(x,z\)。
	\end{solution}
\end{question}

\subsubsection{量词推理规则易错点}

\begin{enumerate}
	\item ES 引入的常量必须是新的，不能在前提或之前推导中已有特殊含义。
	\item UG 只能对“任意个体”推广。若某个常量来自 ES，通常不能对它使用 UG。
	\item \(\exists x\forall yA(x,y)\Rightarrow\forall y\exists xA(x,y)\) 成立，但反向一般不成立。
	\item \(\forall x\exists yA(x,y)\) 一般不能推出 \(\exists y\forall xA(x,y)\)。
\end{enumerate}

\subsection{本章重点定义与定理速查}

\subsubsection{命题逻辑核心定义}

\begin{enumerate}
	\item \textbf{命题}：能够判断真假的陈述句。每个命题有且仅有一个真值。
	\item \textbf{合式公式}：由命题变元、联结词和括号按形成规则有限次构造出的公式。
	\item \textbf{真值指派}：给公式中每个命题变元指定 \(0\) 或 \(1\) 的过程。含 \(n\) 个不同变元时共有 \(2^n\) 种指派。
	\item \textbf{永真式}：任意真值指派下都为真的公式，也称重言式。
	\item \textbf{矛盾式}：任意真值指派下都为假的公式。
	\item \textbf{可满足式}：至少存在一种真值指派使其为真的公式。
	\item \textbf{等值式}：若 \(A\leftrightarrow B\) 是永真式，则称 \(A\) 与 \(B\) 等值，记作 \(A\Leftrightarrow B\)。
	\item \textbf{永真蕴含}：若 \(A\to B\) 是永真式，则称 \(A\) 永真蕴含 \(B\)，记作 \(A\Rightarrow B\)。
\end{enumerate}

\subsubsection{命题逻辑核心定理}

\begin{enumerate}
	\item \textbf{蕴含等值定理}：
	\[
	p\to q\Leftrightarrow \neg p\lor q.
	\]
	考试中化简蕴含式、求范式时通常先用它消去 \(\to\)。
	\item \textbf{逆否定理}：
	\[
	p\to q\Leftrightarrow \neg q\to\neg p.
	\]
	证明题中可把原命题改为逆否命题证明。
	\item \textbf{德摩根律}：
	\[
	\neg(p\land q)\Leftrightarrow \neg p\lor\neg q,\qquad
	\neg(p\lor q)\Leftrightarrow \neg p\land\neg q.
	\]
	多变量情形同理，否定进入括号时要交换 \(\land,\lor\)。
	\item \textbf{分配律}：
	\[
	p\land(q\lor r)\Leftrightarrow(p\land q)\lor(p\land r),
	\]
	\[
	p\lor(q\land r)\Leftrightarrow(p\lor q)\land(p\lor r).
	\]
	它是普通析取范式和合取范式互化的主要工具。
	\item \textbf{主范式存在唯一性定理}：任一含 \(n\) 个命题变元的公式都有唯一的主析取范式和唯一的主合取范式。这里唯一是指在固定变元顺序下，小项或大项编号唯一。
	\item \textbf{小项与大项互补关系}：
	\[
	M_i\Leftrightarrow \neg m_i,\qquad m_i\Leftrightarrow \neg M_i.
	\]
	若公式主析取范式含编号集合 \(S\)，则主合取范式含编号集合 \(\{0,1,\dots,2^n-1\}-S\)。
	\item \textbf{联结词完备定理}：\(\{\neg,\land\}\)、\(\{\neg,\lor\}\)、\(\{\neg,\to\}\)、\(\{\uparrow\}\)、\(\{\downarrow\}\) 都是完备联结词集。
\end{enumerate}

\subsubsection{谓词逻辑核心定义与定理}

\begin{enumerate}
	\item \textbf{个体域}：谓词公式中变量取值的范围。量词命题的真假依赖个体域。
	\item \textbf{谓词}：含个体变元的命题函数，如 \(P(x)\)、\(R(x,y)\)。
	\item \textbf{自由变元}：没有被同名量词约束的变量出现。含自由变元的公式通常不是命题。
	\item \textbf{约束变元}：落在同名量词辖域内的变量出现。
	\item \textbf{闭式}：不含自由变元的谓词公式。闭式具有确定真值。
	\item \textbf{前束范式}：所有量词都位于公式最前端，后面是不含量词的公式。
	\item \textbf{量词否定律}：
	\[
	\neg\forall xP(x)\Leftrightarrow\exists x\neg P(x),
	\]
	\[
	\neg\exists xP(x)\Leftrightarrow\forall x\neg P(x).
	\]
	\item \textbf{有限域展开定理}：若个体域为 \(\{a_1,\dots,a_n\}\)，则
	\[
	\forall xP(x)\Leftrightarrow P(a_1)\land\cdots\land P(a_n),
	\]
	\[
	\exists xP(x)\Leftrightarrow P(a_1)\lor\cdots\lor P(a_n).
	\]
	\item \textbf{量词换名原则}：约束变元可在不与自由变元冲突、不改变辖域的条件下改名。换名是化前束范式前的重要步骤。
\end{enumerate}
